\documentclass[article, 11pt]{IEEEtran}   %Document is an article, 11 pt font, and follows IEEE format

%Aqui van todos los paquetes que van a utilizar para el documento
\usepackage{hyperref}         %Uso paquete hyperref  (hiperlinks)
\usepackage{graphicx}         %Uso paquete  graphicx  (imágenes)
\usepackage{float}            %Uso paquete  float  (posicionar imágenes)
\usepackage{gensymb}          %Uso paquete  gensymb  (simbolos)
\usepackage{amsmath}          %Uso paquete  amsmath  (formulas matemáticas)
\usepackage{tikz}             %Uso paquete  tikz  (crear gráficas)
\usepackage{pgfplots}         %Uso paquete  pgfplots  (crear gráficas)
\usepackage{pgfplotstable}    %Uso paquete  pgfplotstable (regresión lineal)
\usepackage[spanish]{babel}	  %Uso paquete  babel spanish (caracteres en español)
\usepackage{lipsum}  		  %Uso paquete  lipsum (crea dummy text)

\begin{document}

%TITULO DEL REPORTE
%Las palabras inician con mayúscula excepto las preposiciones
\title{Proyecto Final\\ Nombre del Proyecto}


%INFORMACIÓN DEL AUTOR O AUTORES
\author{\IEEEauthorblockN{Nombre del Autor}
\IEEEauthorblockA{\\Universidad Galileo, Teoría de Sistemas II\\
Guatemala, Guatemala}}


% make the title area
\maketitle



%Resumen
%Esta sección es un breve resumen de los aspectos más importantes del proyecto. Se incluye el objetivo del proyecto, los hallazgos más importantes y la conclusión principal.

\begin{abstract}
\lipsum[1-1]
\end{abstract}


%INTRODUCCIÓN
%Esta sección incluye el objetivo del proyecto. Se deben incluir los %siguientes temas:
% Contexto de aprendizaje -> de que trata el proyecto, incluir el background teorico, incluir las referencias!
%El proposito del proyecto -> se incluye la hipótesis que desea probar, o la serie de preguntas que se desean contestar
%Breve resumen de la metodología utilizada y de los resultados que esperan obtener

%Link de ayuda: https://libguides.usc.edu/writingguide/introduction
\section{Introducci\'on}
\lipsum[1-3]


%METODOLOGÍA
%Incluye todos los pasos que se deben seguir para llevar a cabo el experimento. Además debe incluir la justificación de la aplicación de los procedimientos. Se debe incluir cómo se obtuvo la data, qué tipo de data es y cómo fue analizada.
\section{Metodolog\'ia}
\lipsum[1-2]

%RESULTADOS
%Incluye únicamente todos los resultados obtenidos del experimentos. Pueden ser presentados en forma de tabla o gráficas. Incluye pequeños parrafos en donde es explique qué es lo que se esta presentando, o a que resultado poner atención. Sin embargo, no incluyen la interpretación de estos.
\section {Resultados}
\lipsum[1-1]
\subsection{Tabla de datos}
\begin{table}[H]
\centering
\caption{Ejemplo de Tabla con datos}
\label{DataTable1}
\begin{tabular}{|c|c|}
\hline
X [s] & Y[m/s]  \\
\hline
11.7 & 37.8              \\
17.1 & 54     \\  
8.4 & 26.6 \\
6.6 & 21.2 \\
7.4 & 23.4 \\
1.9 & 7   \\
\hline   
\end{tabular}
\end{table}

%EJEMPLO DE COMO UTILIZAR LOS PAQUETES PARA CREAR UNA GRAFICA
\subsection{Gr\'aficas}
\begin{figure}[H]
	\centering
	\begin{tikzpicture}
		\begin{axis}[
   			title=Puntos obtenidos,	%Title of graph
   			xlabel={x[s]},			%label the x-axis
  			ylabel={y[m/s]},		%label the y-axis
   			legend pos=south east,	%put slope legend in bottom right corner
]
			\addplot[blue, only marks] table {CircleData.dat};	%plot points from data table
			\addplot [red] table [y={create col/linear regression={y=y}}]   %best fit line
        {CircleData.dat};
       
			\addlegendentry{slope=$\pgfmathprintnumber{\pgfplotstableregressiona}$}        
        %show slope value
		\end{axis}
	\end{tikzpicture}
	\caption{Ejemplo de gr\'afica}\label{diagram1}   %caption the graph
\end{figure}

%DISCUSIÓN 
%Esta sección incluye la interpretación de los resultados presentados anteriormente. Se debe incluir la explicación, analisis y comparación entre lo que ya se conoce de literatura anterior y los resultados obtenidos. También  en esta sección se deben responder las preguntas propuestas en la introducción y como se prueba o rechaza la hipotesis propuesta.
\section{Discusi\'on}
\lipsum[1-2]

%CONCLUSIÓN
%Esta sección es un resumen de lo aprendido al final de haber desarrollado el proyecto. Es una sintesís de los resultados más importantes y de lo aprendido de la teoría y la práctica.
\section{Conclusi\'on}
\lipsum[1-1]

%FORMATO APA
\begin{thebibliography}{1}

\bibitem{LabWrite}
M.~Carter, E.~Weiebe, and M.~Ferzli, \emph{LabWrite}, \url{https://www.ncsu.edu/labwrite/index_labwrite.htm}, NC~State~University, 2004.

\bibitem{Fullerton}
D.~Fullerton, \emph{APlusPhysics}, \url{http://www.aplusphysics.com/courses/regents/lab_report.html}, 2016.

\bibitem{Svoronos}
T.~Svoronos, \emph{How I Will Write My Dissertation}, \url{http://teddysvoronos.com/2014/12/26/how-i-will-write-my-dissertation-3/}, 2014.

\bibitem{Clarion}
G. Clarion, \emph{Rocklin High School}, 2009.

\end{thebibliography}

\end{document}


